\lysh{15816}{2}
\begin{alist}
\item Είναι $\alpha + \gamma = \ln 2 + \ln 8 = \ln(2\cdot 8) = \ln 16 = \ln 4^2 = 2\ln 4 = 2\beta$.
\item Είναι $\beta + \gamma = \ln 4 + \ln 8 = \ln(4\cdot 8) = \ln 32 = \ln 2^5 = 5\ln 2 = 5\alpha$.
\end{alist}
\textbf{Σημείωση:}
Οι αριθμοί $2, 4, 8$ είναι με αυτή τη σειρά διαδοχικοί όροι γεωμετρικής προόδου αφού $4^2 = 2\cdot 8$, ενώ οι αριθμοί $\alpha = \ln 2, \beta = \ln 4, \gamma = \ln 8$ είναι με αυτή τη σειρά διαδοχικοί όροι αριθμητικής προόδου αφού όπως δείξαμε $2\beta = \alpha + \gamma$. \\
Γενικά αν $\alpha, \beta, \gamma$ είναι με αυτή τη σειρά θετικοί διαδοχικοί όροι γεωμετρικής προόδου, τότε οι αριθμοί $\ln\alpha, \ln\beta, \ln\gamma$ είναι με αυτή τη σειρά διαδοχικοί όροι αριθμητικής προόδου.

